\chapter{Consultas SQL/PostGIS de especial complejidad}
\label{ap:sql}

Este apéndice recoge las tres consultas SQL/PostGIS más complejas del proyecto, más allá de las ya mostradas en la Sección~\ref{sec:enrutamiento} (nodo más cercano y consulta de \texttt{pgr\_dijkstra}). Los parámetros precedidos de dos puntos (\texttt{:nombre}) son los que \texttt{SQLAlchemy} sustituye en tiempo de ejecución, siguiendo la misma convención de los listados del Capítulo~\ref{cap:propuesta}.

% -----------------------------------------------------------------------
% B.1 ÍNDICE DE PELIGROSIDAD Y COSTE SEGURO
% -----------------------------------------------------------------------
\section{Índice de peligrosidad y coste seguro}
\label{ap:sec:indice_peligrosidad}

El Listado~\ref{lst:indice-peligrosidad} calcula, para todos los tramos a la vez, la Ecuación~\ref{eq:coste} (índice de peligrosidad) y la Ecuación~\ref{eq:coste_seguro} (coste seguro), descritas en la Sección~\ref{sec:funcion_coste}. Los topes al percentil 99 (\texttt{:cap\_acc}, \texttt{:cap\_atropello}, \texttt{:cap\_luz}) y los extremos de vulnerabilidad (\texttt{:ivt\_min}, \texttt{:ivt\_max}) se calculan en una consulta previa y se pasan como parámetros.

\begin{lstlisting}[style=SQL, caption={Cálculo del índice de peligrosidad y del coste seguro (\texttt{scripts/07\_calcular\_indice\_seguridad.py}).}, label={lst:indice-peligrosidad}]
UPDATE edges e
SET indice_peligrosidad =
      :w_acc  * LEAST(e.accidentes_100m, :cap_acc) / :cap_acc
    + :w_atr  * LEAST(e.num_atropellos, :cap_atropello) / :cap_atropello
    + :w_luz  * (1 - LEAST(e.farolas_100m, :cap_luz) / :cap_luz)
    + :w_vuln * (v.ivt_agregado - :ivt_min) / (:ivt_max - :ivt_min)
FROM vulnerabilidad_distritos v
WHERE v.cod_distrito = e.cod_distrito;

UPDATE edges
SET cost_seguro = length * (1 + :alpha * indice_peligrosidad),
    reverse_cost_seguro = length * (1 + :alpha * indice_peligrosidad);
\end{lstlisting}

% -----------------------------------------------------------------------
% B.2 DENSIDAD DE FAROLAS POR TRAMO MEDIANTE BUFFER DE PROXIMIDAD
% -----------------------------------------------------------------------
\section{Densidad de farolas por tramo mediante buffer de proximidad}
\label{ap:sec:buffer_farolas}

El Listado~\ref{lst:buffer-farolas} cuenta, para cada tramo, el número de farolas en un radio de 20\,m (\texttt{BUFFER\_M}, Sección~\ref{subsec:iluminacion_datos}) mediante \texttt{ST\_DWithin}, y deriva de ahí la densidad \texttt{farolas\_100m}. El radio se expresa en grados (\texttt{:buffer\_deg}) y no mediante el tipo \texttt{geography}, que sí trabajaría en metros de forma exacta pero no puede aprovechar el índice \texttt{GIST} ya creado sobre la columna \texttt{geometry}: con un radio tan pequeño, la distorsión de trabajar directamente en grados es despreciable frente a la ganancia de poder usar dicho índice sobre una tabla de más de 524\,000 tramos.

\begin{lstlisting}[style=SQL, caption={Densidad de farolas por tramo (\texttt{scripts/03\_load\_farolas.py}).}, label={lst:buffer-farolas}]
UPDATE edges e
SET num_farolas = sub.n,
    farolas_100m = sub.n / (e.length / 100.0)
FROM (
    SELECT e2.id, count(f.id) AS n
    FROM edges e2
    JOIN farolas f
      ON ST_DWithin(f.geom, e2.geom, :buffer_deg)
    GROUP BY e2.id
) sub
WHERE e.id = sub.id AND e.length > 0;
\end{lstlisting}

% -----------------------------------------------------------------------
% B.3 JOIN ESPACIAL TRAMO-DISTRITO CON RESOLUCIÓN DE CASOS LÍMITE
% -----------------------------------------------------------------------
\section{Join espacial tramo-distrito con resolución de casos límite}
\label{ap:sec:join_distrito}

El Listado~\ref{lst:join-distrito} asigna a cada tramo el distrito que contiene su centroide (Sección~\ref{sec:modelado_red}), y resuelve los tramos que quedan sin asignar (el centroide cae fuera de todos los polígonos, típicamente por pequeñas discrepancias entre la red de OSM y los límites administrativos oficiales) asignándolos al distrito más cercano mediante el operador de vecino más cercano (\texttt{<->}) en una subconsulta \texttt{LATERAL}.

\begin{lstlisting}[style=SQL, caption={Join espacial tramo-distrito (\texttt{scripts/06\_load\_distritos.py}).}, label={lst:join-distrito}]
-- Distrito que contiene el centroide del tramo.
UPDATE edges e
SET cod_distrito = sub.cod_distrito
FROM (
    SELECT e2.id, d.cod_distrito
    FROM edges e2
    JOIN distritos d ON ST_Within(ST_Centroid(e2.geom), d.geom)
) sub
WHERE e.id = sub.id;

-- Tramos sin distrito asignado: al distrito más cercano.
UPDATE edges e
SET cod_distrito = sub.cod_distrito
FROM (
    SELECT e2.id, d.cod_distrito
    FROM edges e2
    CROSS JOIN LATERAL (
        SELECT cod_distrito
        FROM distritos
        ORDER BY geom <-> ST_Centroid(e2.geom)
        LIMIT 1
    ) d
    WHERE e2.cod_distrito IS NULL
) sub
WHERE e.id = sub.id;
\end{lstlisting}
